\documentclass[10pt,a4paper]{article}
\usepackage{fontspec,amsmath,amssymb,graphicx}
\usepackage[margin=18mm]{geometry}
\setmainfont{Times New Roman}
\setlength{\parindent}{0pt}
\setlength{\parskip}{5pt}
\allowdisplaybreaks
\sloppy
\emergencystretch=2em
\begin{document}
\section*{BT\_3\_1\_3\_03 — Lời giải bổ sung}
\subsection*{Kiến thức và trình tự giải}
\textbf{Bước 1 — Nhận dạng dữ liệu và tốc độ.} Viết $u_{tt}-a^2u_{xx}=F$, $u(x,0)=f(x)$, $u_t(x,0)=g(x)$; tốc độ là $a>0$, không phải $a^2$. Với dữ liệu trơn, cần $f\in C^2$, $g\in C^1$. Nếu dữ liệu có góc hoặc bước nhảy, công thức vẫn mô tả sóng từng miền nhưng phải xác định nơi mất tính trơn.

\textbf{Bước 2 — Tìm hai sóng chạy.} Đặt $\xi=x-at$, $\eta=x+at$. Quy tắc dây chuyền cho $u_{tt}-a^2u_{xx}=-4a^2u_{\xi\eta}$, nên khi $F=0$, $u=A(\xi)+B(\eta)$. Tại $t=0$, $A+B=f$, $a(B'-A')=g$. Chọn một nguyên hàm $H'=g$ rồi lấy $A=f/2-H/(2a)$, $B=f/2+H/(2a)$; hằng số của $H$ sẽ triệt tiêu.


\[
u_0(x,t)=\frac{f(x-at)+f(x+at)}2+\frac{H(x+at)-H(x-at)}{2a}.
\]


\textbf{Bước 3 — Tính tích phân theo dữ liệu thật.} Tìm nguyên hàm trên từng khoảng, nối liên tục nếu $g$ có bước nhảy; thay hai cận $x\pm at$. Dùng công thức cộng lượng giác trước khi kết luận chu kỳ. Với dữ liệu chỉ thị, tích phân bằng độ dài giao của $[x-at,x+at]$ với đoạn hỗ trợ, nhân hệ số của $g$.

\textbf{Bước 4 — Nếu có nguồn, thêm Duhamel.} Nghiệm riêng có dữ liệu đầu bằng không là


\[
u_F(x,t)=\frac1{2a}\int_0^t\int_{x-a(t-s)}^{x+a(t-s)}F(y,s)\,dy\,ds.
\]


Tính tích phân trong trước. Nếu $F$ là đa thức hoặc một mode lượng giác, có thể đoán $u_F$ rồi giải ODE theo $t$, nhưng vẫn phải buộc $u_F(x,0)=(u_F)_t(x,0)=0$. Nghiệm cuối là $u_0+u_F$.

\textbf{Bước 5 — Kiểm tra và đọc ý nghĩa.} Đạo hàm công thức: $u_x=(f'(\xi)+f'(\eta))/2+(g(\eta)-g(\xi))/(2a)$, $u_t=a(f'(\eta)-f'(\xi))/2+(g(\eta)+g(\xi))/2$. Tại $t=0$ thu lại $f,g$; đạo hàm tiếp cho PDE. Các điểm gãy $s=b$ phát sinh đường $x\pm at=b$. Chu kỳ phải đúng cho mọi $x,t$, không chỉ sau khi một mặt sóng đi qua một điểm cố định.
\subsection*{Lời giải áp dụng cho bài}
\textbf{Bước 1. Tách sóng li độ và vận tốc.} $a=5$, $f=10\chi_{[-1,1]}$, $g=\chi_{[-1,1]}$. Tại $x=15$, cận phải $15+5t$ luôn ngoài hỗ trợ. Cận trái $15-5t$ nằm trong đoạn khi $14/5\le t\le16/5$. Đặt $\ell(t)=|[15-5t,15+5t]\cap[-1,1]|$.

\textbf{Bước 2. Viết áp suất theo từng giai đoạn.}
\[u(15,t)=\begin{cases}0&0\le t<14/5,\\5+(5t-14)/10&14/5\le t\le16/5,\\1/5&t>16/5.\end{cases}\]
Hạng $5$ là một nửa của li độ $10$; hạng thứ hai là $\ell/(2a)$. Ở hai đầu mút, giá trị phụ thuộc quy ước chỉ thị đóng; các vết một phía vẫn như trên.

\textbf{Bước 3. So sánh ngưỡng.} Trong khoảng giữa, hàm tăng, lớn nhất $5+2/10=26/5=5.2$, thấp hơn ngưỡng $8$. Sau xung chỉ còn $1/5$. Do đó áp suất không vượt ngưỡng. Không thể chỉ xét tích phân vận tốc mà bỏ sóng li độ gây bước nhảy $5$.
\begin{center}\includegraphics[width=.97\linewidth,height=11cm,keepaspectratio]{Do_thi/BT_3_1_3_03.png}\end{center}
\textit{Đồ thị minh họa; với nghiệm số, xem phương pháp và tham số trong mã Python kèm theo.}\end{document}